\section{Introduction}
\label{sec:introduction}

Mechanized cryptography needs several descriptions of one algorithm.  A proof
of correctness observes a probability distribution; an implementation-level
argument records which operations were executed; a complexity proof supplies a
worst-case bound; and a security definition admits only adversaries carrying an
appropriate polynomial-time certificate.  If these descriptions are defined
independently, they can drift.  In particular, a costed interpreter can disagree
with the probability game, a ``runtime'' field can become metadata unrelated to
the run it purports to bound, and an interactive security definition can mention
protocols or corruption policies that never enter the actual execution.

This paper describes the current architecture of \LeanCrypto, an experimental
cryptographic verification library implemented in Lean~4
\cite{deMouraUllrich2021Lean4} and built on mathlib
\cite{mathlib2020}.  The architecture is intended to support game-based and
universally composable security in a shared executable substrate.  Its central
invariant is deliberately stronger than a naming convention.

\begin{designprinciple}
Every public probability semantics is obtained by erasing cost from one
authoritative exact execution.  Every exact, measured, and polynomial
complexity result is a proposition or certificate indexed by that same
execution.
\end{designprinciple}

Four separations make the invariant practical.  First, exact execution and
cost-erased probability are different views of one joint value--cost
distribution.  Second, the cost produced on a path and a proof that it is below
a budget are different objects.  Third, an abstract resource---which may be
multidimensional or order-sensitive---is distinct from the natural-number
observation used by the existing asymptotic layer.  Fourth, a security game is
distinct from the certificate that admits a machine as a PPT adversary.  These
separations run from primitive operations through typed programs and oracles to
the interactive-machine kernel.

ElGamal encryption is a representative example.  Its encryption algorithm is
a typed program whose exact handler samples one scalar, performs two scalar
actions, and performs one group operation.  Correctness consumes the erased
distribution of that program.  Efficiency attaches an input-dependent bound
and a measured runtime certificate to the same program; it does not restate the
algorithm in a second AST.  The one-time pad, discrete logarithm, and
Decisional Diffie--Hellman examples follow the same boundary.

\paragraph{Contributions.}
The present implementation provides:

\begin{itemize}
  \item a strict import hierarchy separating security parameters,
    probability, exact cost semantics, asymptotics, machine certificates,
    game-based notions, and UC execution;
  \item exact randomized resource semantics over an ordered, potentially
    noncommutative additive monoid, with a single support-wise path-bound
    predicate and a monotone additive observation into \(\Nat\);
  \item typed result-indexed primitive signatures whose exact handlers,
    cost-erased mathematical laws, and upper bounds are separate records;
  \item a typed program language with one interpreter, an exact structural
    execution relation, a support-equivalence theorem, and compositional
    input-dependent bounds;
  \item a unified dependent machine hierarchy and a closed exact-to-measured-
    to-polynomial certificate chain, including a single oracle interpreter and
    certified caller--oracle composition; and
  \item a typed FIFO UC execution layer with capabilities, adaptive
    corruption, explicit kernel charging, certified finite fuel, real/ideal
    wiring, and a context-composition theorem derived from operational
    simulation.
\end{itemize}

The goal of the paper is architectural: to state what the current Lean
interfaces guarantee, how the certificates connect, and which obligations
remain with an instantiation.  It does not claim a new cryptographic reduction
or a physical CPU-time model.  The distinction is made explicit in
\Cref{sec:boundaries}.

\paragraph{Positioning.}
The library shares the broader aim of foundational mechanized cryptography
exemplified by SSProve \cite{CSF:AHRVWH21}.  Its interactive layer follows the
real/ideal quantifier discipline of universal composability
\cite{FOCS:Canetti01}, while its engineering focus differs from systems such as
EasyUC \cite{CSF:CanStoVar19}: the contribution described here is a single
Lean-native exact resource and certificate hierarchy spanning typed programs,
adaptive oracles, and an executable UC kernel.  A systematic feature or proof
engineering comparison is left for later work.

\paragraph{Organization.}
\Cref{sec:hierarchy} presents the dependency discipline.
\Cref{sec:cost-algebra-program} develops exact resources, typed algebras, and programs.
\Cref{sec:machines-oracles} gives the machine and oracle certificate chains.
\Cref{sec:uc} describes the interactive execution stack.
\Cref{sec:case-studies,sec:evaluation} summarize concrete instantiations and
mechanized validation, followed by soundness boundaries and future work.
